\begin{afterword}
\summaryhead

The post opens with a fable. A tribe's wise man, asked why water falls from the sky, says that sky spirits battle and their blood drips down; asked where the sky spirits come from, he whispers, ``From the before time.''

Faced with something you cannot explain, the author says, you can Ignore it, try to Explain it, or Worship it, which means enjoying the sensation of mysteriousness. Each explanation raises the same choice again: life is explained by chemistry, chemistry by atoms, atoms by nuclei, nuclei by quarks, and the quarks by the Big Bang. It is tempting to conclude that the chain must end in an ``Explanation That Needs No Explanation,'' the only explanation worth knowing. The author then points out that this conclusion was itself Worship. The Explain response would have been to find the regress paradoxical and wonder how it is resolved. Treating the question as unimportant, or putting it off, is Ignore. The post ends: ``Select your option wisely.''

\responsehead

The post is a fable, a three-way classification and one twist: the reader follows an argument to a final explanation and is then told that the argument was Worship. The distinction behind it is clear and useful. A question can be treated as open, or as a sacred mystery, or dropped. The problems are in what the post leaves unsaid.

Much of the post restates the two posts of the week before. ``Semantic Stopsigns'' (24 August 2007) opens with the same shape of dialogue, a chain of questions ending at a mythic origin, and already traces the regress back to the Big Bang and the laws of physics. ``Mysterious Answers to Mysterious Questions'' (25 August 2007) already calls the enjoyment of mystery worship. What this post adds is the three commands and the twist.

The twist does not say what it condemns. The paragraph labelled Worship contains a claim, that the chain of explanations must end somewhere, and a valuation, that the end is the only explanation worth knowing. The first is the old question of a first cause. The post does not say whether concluding that the chain ends is Worship, or only the reverence that comes with it.

The post then stops. Its Explain response to the regress is to wonder how the paradox is resolved. ``Semantic Stopsigns'' had also declined to discuss it. The last line tells the reader to choose wisely among the three options, and the post has given no way to tell when setting a question aside is the wise choice.

In short: a clear distinction between treating a question as open and treating it as sacred, made mostly with material from the week before, with a twist that does not say which part of the reasoning was Worship, and a closing choice it gives no way to make.
\end{afterword}
